\documentclass[10pt]{article}

\usepackage[utf8]{inputenc}
\usepackage[T1]{fontenc}
\usepackage{amsmath,amssymb,mathrsfs}
\usepackage{microtype}
\usepackage[margin=1in]{geometry}

\title{Exponential Sums with Monomials}
\author{Étienne Fouvry\\
  \small Université de Paris-Sud, 91405 Orsay, France
  \and
  Henryk Iwaniec\\
  \small Rutgers University, New Brunswick, New Jersey 08903}
\date{Communicated by H. L. Montgomery\\Received April 6, 1988}

\begin{document}

\maketitle
\begin{center}
  \small\textsc{Journal of Number Theory} \textbf{33} (1989), 311--333
\end{center}

\section*{1. Introduction}
The exponential sums of type


\begin{equation*}
\sum_{m} e(f(m)), \tag{1.1}
\end{equation*}


where $m$ ranges over integers from an interval and $f(m)$ is a smooth function, play a central role in analytic number theory. General methods of estimating such sums were established by H. Weyl [9], J. G. van der Corput [2], and I. M. Vinogradov [8]. Later, the exponential sums in several variables


\begin{equation*}
\sum_{m_{1}} \cdots \sum_{m_{j}} e\left(f\left(m_{1}, \ldots, m_{j}\right)\right) \tag{1.2}
\end{equation*}


were introduced to enhance these methods as well as for their own importance. Our main interest in this paper is to estimate the exponential sums in which $f$ is a monomial function,


\begin{equation*}
f\left(m_{1}, \ldots, m_{j}\right)=x m_{1}^{\alpha_{1}} \cdots m_{j}^{\alpha_{j}} . \tag{1.3}
\end{equation*}


We regard such sums as a special case of bilinear forms


\begin{equation*}
\mathscr{B}_{\varphi \psi}(\mathscr{X}, \mathscr{Y})=\sum_{r} \sum_{s} \varphi_{r} \psi_{s} e\left(x_{r} y_{s}\right), \tag{1.4}
\end{equation*}


where $\mathscr{X}=\left(x_{r}\right), \mathscr{Y}=\left(y_{s}\right)$ are finite sequences of real numbers with


\begin{equation*}
\left|x_{r}\right| \leqslant X, \quad\left|y_{s}\right| \leqslant Y, \tag{1.5}
\end{equation*}


say, and $\varphi_{r}, \psi_{s}$ are complex numbers. Our arguments are based on the following general inequality of [1].

Proposition 1. We have


\begin{equation*}
\left|\mathscr{B}_{\varphi \psi}(\mathscr{X}, \mathscr{Y})\right|^{2} \leqslant 20(1+X Y) \mathscr{B}_{\varphi}(\mathscr{X}, Y) \mathscr{B}_{\psi}(\mathscr{Y}, X) \tag{1.6}
\end{equation*}


with

$$
\mathscr{B}_{\varphi}(\mathscr{X}, Y)=\sum_{\left|x_{r_{1}}-x_{r_{2}}\right| \leqslant Y^{-1}}\left|\varphi_{r_{1}} \varphi_{r_{2}}\right|
$$

and $\mathscr{B}_{\psi}(\mathscr{Y}, X)$ defined similarly.\\
We begin by applying (1.6) directly. Next we introduce a number of innovations to combine with (1.6), giving deeper results. Finally, we illustrate how to use these results for estimating sums which occur in sieve problems for short intervals.

Notation and conventions:

$$
\begin{aligned}
& {[x]=\max \{k \in \mathbb{Z} ; k \leqslant x\} .} \\
& \|x\|=\min \{|x-k| ; k \in \mathbb{Z}\} . \\
& e(z)=\exp (2 \pi i z) .
\end{aligned}
$$

$f \ll g$ means $|f| \leqslant c g$ with some positive constant $c$.\\
$f=O(g)$ means $f \ll g$.\\
$f \asymp g$ means $c_{1} f<g<c_{2} g$ with some positive, unspecified constants $c_{1}, c_{2}$.

The constants implied in the symbols $O$ and $\ll$ may depend, without mentioning, on those implied in the relevant relations $f \asymp g$.\\
$\|\varphi\|$ stands for the $\ell_{2}$-norm of the sequence $\varphi=\left(\varphi_{r}\right)$, i.e.,

$$
\|\varphi\|=\left(\sum_{r}\left|\varphi_{r}\right|^{2}\right)^{1 / 2} .
$$

The symbol $\blacksquare$ indicates the end of a proof.

\section*{2. Direct Results}
By Proposition 1 we immediately obtain the following.

Corollary 1. Suppose that the sequences $\mathscr{X}$ and $\mathscr{Y}$ are $A$-spaced and $B$-spaced, respectively; that is, $\left|x_{r_{1}}-x_{r_{2}}\right| \geqslant A$ and $\left|y_{s_{1}}-y_{s_{2}}\right| \geqslant B$ for $r_{1} \neq r_{2}$ and $s_{1} \neq s_{2}$, respectively. We then have

$$
\left|\mathscr{B}_{\varphi \psi}(\mathscr{X}, \mathscr{Y})\right| \leqslant 5(1+X Y)^{1 / 2}\left(1+\frac{1}{A Y}\right)^{1 / 2}\left(1+\frac{1}{B X}\right)^{1 / 2}\|\varphi\|\|\psi\| .
$$

Yet this result is ineffective since the spacing problem is not resolved. We easily handle that in the following situation.

Theorem 1. Let $f$ and $g$ be smooth functions such that

$$
\begin{array}{ll}
f \asymp F, & f^{\prime} \asymp F M^{-1}, \\
g \asymp G, & g^{\prime} \asymp G N^{-1},
\end{array}
$$

with $F, G, M, N$ positive, and let $\varphi_{m}, \psi_{n}$ be complex numbers. We then have

$$
\begin{aligned}
S_{\varphi \psi}(M, N) & =\sum_{m \sim M} \sum_{n \sim N} \varphi_{m} \psi_{n} e(f(m) g(n)) \\
& \ll(F G)^{-1 / 2}(F G+M)^{1 / 2}(F G+N)^{1 / 2}\|\varphi\|\|\psi\| .
\end{aligned}
$$

Proof. It follows from Corollary 1 and from the relations

$$
\left|f\left(m_{1}\right)-f\left(m_{2}\right)\right| \asymp\left|m_{1}-m_{2}\right| M^{-1} F, \quad\left|g\left(n_{1}\right)-g\left(n_{2}\right)\right| \asymp\left|n_{1}-n_{2}\right| N^{-1} G .
$$ \(\square\)

Next we investigate the spacing of binomial points $m^{\alpha} n^{\beta}$.\\
Lemma 1. Let $\alpha \beta \neq 0, \Delta>0, M \geqslant 1$ and $N \geqslant 1$. Let $\mathscr{A}(M, N ; \Delta)$ be the number of quadruples ( $m, \tilde{m}, n, \tilde{n}$ ) such that

$$
\left|\left(\frac{\tilde{m}}{m}\right)^{\alpha}-\left(\frac{\tilde{n}}{n}\right)^{\beta}\right|<\Delta,
$$

with $M \leqslant m, \tilde{m}<2 M$ and $N \leqslant n, \tilde{n}<2 N$. We then have

$$
\mathscr{A}(M, N ; \Delta) \ll M N \log 2 M N+\Delta M^{2} N^{2} .
$$

Proof. We divide the solutions into classes each one having fixed values $(m, \tilde{m})=\mu$ and $(n, \tilde{n})=\nu$, say. In a given class the points $(\tilde{m} / m)^{\alpha}$ are spaced by $c(\alpha) \mu^{2} M^{-2}$ and the points $(\tilde{n} / n)^{\beta}$ are spaced by $c(\beta) \nu^{2} N^{-2}$, where $c(\alpha)$\\
and $c(\beta)$ are positive constants. Therefore, by the Dirichlet box principle each class contains

$$
\begin{aligned}
& \ll \min \left\{\mu^{-2} M^{2}\left(1+\Delta \nu^{-2} N^{2}\right), \nu^{-2} N^{2}\left(1+\Delta \mu^{-2} M^{2}\right)\right\} \\
& =\min \left\{\mu^{-2} M^{2}, \nu^{-2} N^{2}\right\}+\Delta \mu^{-2} \nu^{-2} M^{2} N^{2}
\end{aligned}
$$

points. Summing over $\mu$ and $\nu$ we complete the proof. \(\blacksquare\)\\
Proposition 1 and Lemma 1 yield an estimate for the exponential sum (1.2) in which $f$ is a quadrinomial function.

Theorem 2. Let $\alpha_{j} \neq 0$ and $M_{j} \geqslant 1$ for $j=1,2,3,4$, let $x>0$, and let $\varphi_{m_{1}m_{2}}, \psi_{m_{3}m_{4}}$ be complex numbers with $\left|\varphi_{m_{1}m_{2}}\right| \leqslant 1$ and $\left|\psi_{m_{3}m_{4}}\right| \leqslant 1$. We then have

$$
\begin{aligned}
S_{\varphi \psi}\left(M_{1}, M_{2}, M_{3}, M_{4}\right)= & \sum_{\substack{m_{j} \sim M_{j}\\1\leqslant j\leqslant4}} \varphi_{m_{1}m_{2}} \psi_{m_{3}m_{4}} e\left(x \frac{m_{1}^{\alpha_{1}} m_{2}^{\alpha_{2}} m_{3}^{\alpha_{3}} m_{4}^{\alpha_{4}}}{M_{1}^{\alpha_{1}} M_{2}^{\alpha_{2}} M_{3}^{\alpha_{3}} M_{4}^{\alpha_{4}}}\right) \\
\ll & \left\{\left(x M_{1} M_{2} M_{3} M_{4}\right)^{1 / 2}+M_{1} M_{2}\left(M_{3} M_{4}\right)^{1 / 2}\right. \\
& \left.+\left(M_{1} M_{2}\right)^{1 / 2} M_{3} M_{4}+x^{-1 / 2} M_{1} M_{2} M_{3} M_{4}\right\} \\
& \times \log\left(2 M_{1} M_{2} M_{3} M_{4}\right) .
\end{aligned}
$$

Proof. We apply (1.6) to the two sequences
\[
\mathscr{X}=\left(x m_{1}^{\alpha_{1}} m_{2}^{\alpha_{2}} M_{1}^{-\alpha_{1}} M_{2}^{-\alpha_{2}}\right),
\qquad
\mathscr{Y}=\left(m_{3}^{\alpha_{3}} m_{4}^{\alpha_{4}} M_{3}^{-\alpha_{3}} M_{4}^{-\alpha_{4}}\right),
\]
obtaining

$$
S_{\varphi \psi} \ll\left((1+x) S_{\varphi} S_{\psi}\right)^{1 / 2},
$$

where

$$
S_{\varphi}=\#\left\{\left(m_{1}, m_{2}, \tilde{m}_{1}, \tilde{m}_{2}\right) ;\left|m_{1}^{\alpha_{1}} m_{2}^{\alpha_{2}}-\tilde{m}_{1}^{\alpha_{1}} \tilde{m}_{2}^{\alpha_{2}}\right| \ll x^{-1} M_{1}^{\alpha_{1}} M_{2}^{\alpha_{2}}\right\}
$$

and $S_{\psi}$ is defined similarly. Hence $S_{\varphi}$ is bounded by the number of solutions to

$$
\left|\left(\frac{\tilde{m}_{1}}{m_{1}}\right)^{\alpha_{1}}-\left(\frac{\tilde{m}_{2}}{m_{2}}\right)^{\alpha_{2}}\right| \ll \frac{1}{x},
$$

so, by Lemma 1 we have

$$
S_{\varphi} \ll M_{1} M_{2} \log\left(2 M_{1} M_{2}\right)+x^{-1} M_{1}^{2} M_{2}^{2} .
$$

Similarly

$$
S_{\psi} \ll M_{3} M_{4} \log\left(2 M_{3} M_{4}\right)+x^{-1} M_{3}^{2} M_{4}^{2} .
$$

These estimates and (1.6) yield the assertion of Theorem 2 provided $x \geqslant 1$. If $x<1$ the assertion is trivial.

Remarks. Theorem 2 is useful in the range $1 \ll x \ll M_{1} M_{2} M_{3} M_{4}$. The sharpest bound is attained for $M_{1} M_{2} M_{3} M_{4} \sim x^{2}$. The reason that Theorem 2 is trivial for $x \ll 1$ is that the exponential factor does not oscillate and the reason that it is also trivial for $x \gg M_{1} M_{2} M_{3} M_{4}$ is that the exponential factor oscillates too rapidly. In the latter case the method fails because the number of terms is too small to control the oscillations. In order to overcome this problem, in the next section, we appeal to Weyl's method which creates additional points of summation and at the same time it reduces the oscillatory behaviour of the exponential factor. There is a double price for this operation. First is that Weyl's method shifts the argument of the exponential function destroying its monomial character and consequently it causes serious problems about the spacing of the resulting points. Second is the use of Cauchy's inequality which halves the final saving.

\section*{3. Two Combinations with the Weyl Shift}
Weyl's method depends on the following inequality.
Lemma 2. Let $L>K, Q>0$, and $z_{k}$ be complex numbers. We then have

$$
\left|\sum_{K \leqslant k<L} z_{k}\right|^{2} \leqslant\left(2+\frac{L-K}{Q}\right) \sum_{|q|<Q}\left(1-\frac{|q|}{Q}\right) \sum_{K \leqslant k-q, k+q<L} z_{k-q} \bar{z}_{k+q} .
$$

Proof. It is similar to that of Lemma 5.10 in [6]. \(\square\)

Our aim is to prove the following.
Theorem 3. Let $\alpha, \alpha_{1}, \alpha_{2}$ be real constants such that $\alpha \neq 1$ and $\alpha \alpha_{1} \alpha_{2} \neq 0$. Let $M, M_{1}, M_{2}, x \geqslant 1$, and $\varphi_{m}, \psi_{m_{1} m_{2}}$ be complex numbers with $\left|\varphi_{m}\right| \leqslant 1$ and $\left|\psi_{m_{1} m_{2}}\right| \leqslant 1$. We then have

$$
\begin{aligned}
S_{\varphi \psi}\left(M, M_{1}, M_{2}\right)= & \sum_{m \sim M} \sum_{m_{1} \sim M_{1}} \sum_{m_{2} \sim M_{2}} \varphi_{m} \psi_{m_{1}m_{2}} e\left(x \frac{m^{\alpha} m_{1}^{\alpha_{1}} m_{2}^{\alpha_{2}}}{M^{\alpha} M_{1}^{\alpha_{1}} M_{2}^{\alpha_{2}}}\right) \\
\ll & \left\{x^{1 / 4} M^{1 / 2}\left(M_{1} M_{2}\right)^{3 / 4}+M^{7 / 10} M_{1} M_{2}\right. \\
& \left.+M\left(M_{1} M_{2}\right)^{3 / 4}+x^{-1 / 4} M^{11 / 10} M_{1} M_{2}\right\}\left(\log\left(2 M M_{1} M_{2}\right)\right)^{2} .
\end{aligned}
$$

Proof. By Lemma 2 we get that

$$
\left|S_{\varphi \psi}\right|^{2} \ll Q^{-1} M M_{1} M_{2}\left(M M_{1} M_{2}+\left|S\left(Q_{0}\right)\right| \log 2 Q\right)
$$

for any $Q \leqslant \frac{1}{3} M$ and some $Q_{0} \leqslant Q$, where

$$
S\left(Q_{0}\right)=\sum_{q \sim Q_{0}}\left(1-q Q^{-1}\right) \sum_{\substack{m+q\sim M\\m-q\sim M}} \sum_{m_{1}\sim M_{1}} \sum_{m_{2}\sim M_{2}} \varphi_{m+q} \overline{\varphi}_{m-q} e\left(x \frac{t(m, q) m_{1}^{\alpha_{1}} m_{2}^{\alpha_{2}}}{M^{\alpha} M_{1}^{\alpha_{1}} M_{2}^{\alpha_{2}}}\right)
$$

and

$$
t(m, q)=(m+q)^{\alpha}-(m-q)^{\alpha} \asymp Q_{0} M^{\alpha-1}
$$

By (1.6) we obtain


\begin{equation*}
S\left(Q_{0}\right)\ll\left(\mathscr{A} \mathscr{B} x Q_{0} M^{-1}\right)^{1 / 2} . \tag{3.1}
\end{equation*}


Here $\mathscr{A}$ is the number of quadruples $\left(m_{1}, m_{2}, \tilde{m}_{1}, \tilde{m}_{2}\right)$ such that

$$
\left|\left(\frac{\tilde{m}_{1}}{m_{1}}\right)^{\alpha_{1}}-\left(\frac{\tilde{m}_{2}}{m_{2}}\right)^{\alpha_{2}}\right|<\left(x Q_{0}\right)^{-1} M,
$$

thus by Lemma 1 we have


\begin{equation*}
\mathscr{A} \ll M_{1} M_{2} \log\left(2 M_{1} M_{2}\right)+\left(x Q_{0}\right)^{-1} M M_{1}^{2} M_{2}^{2} \tag{3.2}
\end{equation*}


It remains to estimate $\mathscr{B}$, which stands for the number of quadruples ( $m, \tilde{m}, q, \tilde{q}$ ) such that

$$
|t(m, q)-t(\tilde{m}, \tilde{q})|<x^{-1} M^{\alpha} .
$$

We shall consider this problem in the next section. From Proposition 2 we obtain


\begin{equation*}
\mathscr{B} \ll Q_{0} M\left(1+x^{-1} M^{2}\right)(\log(2 M))^{4} \tag{3.3}
\end{equation*}


provided $3 Q_{0} \leqslant M^{3 / 5}$. Combining (3.1), (3.2), and (3.3) we infer

$$
\begin{aligned}
S\left(Q_{0}\right) \ll & Q_{0}\left(x M_{1} M_{2}\right)^{1 / 2}\left(1+x^{-1} Q_{0}^{-1} M M_{1} M_{2}\right)^{1 / 2} \\
& \times\left(1+x^{-1} M^{2}\right)^{1 / 2}\left(\log\left(2 M M_{1} M_{2}\right)\right)^{5 / 2}
\end{aligned}
$$

The worst value of $Q_{0}$ is $Q_{0}=Q$. Setting $Q=\frac{1}{3} M^{3 / 5}$ we conclude the proof. \(\blacksquare\)

Next we estimate the bilinear forms of type

$$
S_{\varphi \psi}\left(M_{1}, M_{2}\right)=\sum_{m_{1} \sim M_{1}} \sum_{m_{2} \sim M_{2}} \varphi_{m_{1}} \psi_{m_{2}} e\left(x \frac{m_{1}^{\alpha_{1}} m_{2}^{\alpha_{2}}}{M_{1}^{\alpha_{1}} M_{2}^{\alpha_{2}}}\right)
$$

by alternate application of the Weyl shift.

Theorem 4. Let $\alpha_{1}, \alpha_{2}$ be real constants different from 0 and 1. Let $M_{1}, M_{2}, x \geqslant 1$, and $\varphi_{m_{1}}, \psi_{m_{2}}$ be complex numbers with $\left|\varphi_{m_{1}}\right| \leqslant 1$ and $\left|\psi_{m_{2}}\right| \leqslant 1$. We then have

$$
\begin{aligned}
S_{\varphi \psi}\left(M_{1}, M_{2}\right)\ll & \left\{x^{1 / 8}\left(M_{1} M_{2}\right)^{3 / 4}+M_{1}^{4 / 5} M_{2}+M_{1} M_{2}^{17 / 20}\right. \\
& \left.+x^{-1 / 8}\left(M_{1} M_{2}\right)^{21 / 20}\right\}\left(\log\left(2 M_{1} M_{2}\right)\right)^{2} .
\end{aligned}
$$

Proof. By Lemma 2 we get that

$$
\begin{aligned}
& \left|S_{\varphi \psi}\right|^{2} \ll Q_{1}^{-1} M_{1}^{2} M_{2}^{2}+Q_{1}^{-1} M_{1} M_{2}\log(2 M_{1}) \\
& \quad\times\left|\sum_{q_{1} \sim Q_{1}^{*}}\left(1-\frac{q_{1}}{Q_{1}^{*}}\right) \sum_{m_{1} \sim M_{1}} \sum_{m_{2} \sim M_{2}} \varphi_{m_{1}+q_{1}} \overline{\varphi}_{m_{1}-q_{1}} e\left(x t\left(m_{1}, q_{1}\right) m_{2}^{\alpha_{2}} M_{1}^{-\alpha_{1}} M_{2}^{-\alpha_{2}}\right)\right|
\end{aligned}
$$

for any $Q_{1} \leqslant \frac{1}{3} M_{1}$ and some $Q_{1}^{*} \leqslant Q_{1}$. Another application of Lemma 2 gives

$$
\left|S_{\varphi \psi}\right|^{4} \ll Q_{1}^{-2} M_{1}^{4} M_{2}^{4}+Q_{1}^{-2} Q_{2}^{-1} Q_{1}^{*} M_{1}^{3} M_{2}^{3}\left(Q_{1}^{*} M_{1} M_{2}+S\right)\left(\log\left(2 M_{1} M_{2}\right)\right)^{3},
$$

where

$$
S \leqslant \sum_{\substack{m_{1} \sim M_{1} \\ m_{2} \sim M_{2}}} \sum_{\substack{q_{1} \sim Q_{1}^{*} \\ q_{2} \sim Q_{2}^{*}}}\left(1-\frac{q_{2}}{Q_{2}}\right) e\left(x t\left(m_{1}, q_{1}\right) t\left(m_{2}, q_{2}\right) M_{1}^{-\alpha_{1}} M_{2}^{-\alpha_{2}}\right)
$$

for any $Q_{2} \leqslant \frac{1}{3} M_{2}$ and some $Q_{2}^{*} \leqslant Q_{2}$. By Proposition 1 we obtain

$$
S \ll\left(\mathscr{B}_{1} \mathscr{B}_{2} x Q_{1}^{*} Q_{2}^{*} M_{1}^{-1} M_{2}^{-1}\right)^{1 / 2},
$$

where $\mathscr{B}_{1}$ is the number of solutions to the inequality

$$
\left|t\left(m_{1}, q_{1}\right)-t\left(\tilde{m}_{1}, \tilde{q}_{1}\right)\right| \ll\left(x Q_{2}^{*}\right)^{-1} M_{1}^{\alpha_{1}} M_{2}
$$

in integers $m_{1}, \tilde{m}_{1} \sim M_{1}$ and $q_{1}, \tilde{q}_{1} \sim Q_{1}^{*}$ and $\mathscr{B}_{2}$ is defined similarly. Assuming that $3 Q_{1} \leqslant M_{1}^{3 / 5}$ and $3 Q_{2} \leqslant M_{2}^{3 / 5}$ we are entitled to use Proposition 2 giving

$$
\mathscr{B}_{1} \ll Q_{1}^{*} M_{1}\left(1+\frac{M_{1}^{2} M_{2}}{x Q_{2}^{*}}\right)\left(\log(2 M_{1})\right)^{4} .
$$

A similar inequality holds for $\mathscr{B}_{2}$. From both results we get

$$
S \ll x^{1 / 2} Q_{1}^{*} Q_{2}^{*}\left(1+\frac{M_{1}^{2} M_{2}}{x Q_{2}^{*}}\right)^{1 / 2}\left(1+\frac{M_{1} M_{2}^{2}}{x Q_{1}^{*}}\right)^{1 / 2}\left(\log\left(2 M_{1} M_{2}\right)\right)^{4} .
$$

Now observe that the worst values for $Q_{1}^{*}, Q_{2}^{*}$ are $Q_{1}^{*}=Q_{1}$ and $Q_{2}^{*}=Q_{2}$. Setting $Q_{1}=\frac{1}{3} M_{1}^{3 / 5}$ and $Q_{2}=\frac{1}{3} M_{2}^{3 / 5}$ the resulting inequalities finally yield

$$
\begin{aligned}
S_{\varphi \psi} \ll & M_{1} M_{2}\left(M_{1}^{-3 / 10}+M_{2}^{-3 / 20}\right) \log\left(2 M_{1} M_{2}\right) \\
& +x^{1 / 8}\left(M_{1} M_{2}\right)^{3 / 4}\left(1+x^{-1} M_{1}^{2} M_{2}^{2 / 5}\right)^{1 / 8} \\
& \times\left(1+x^{-1} M_{1}^{2 / 5} M_{2}^{2}\right)^{1 / 8}\left(\log\left(2 M_{1} M_{2}\right)\right)^{2} \\
& \ll\left(M_{1} M_{2}^{17 / 20}+M_{1}^{4 / 5} M_{2}+x^{1 / 8} M_{1}^{3 / 4} M_{2}^{3 / 4}+x^{-1 / 8} M_{1}^{21 / 20} M_{2}^{21 / 20}\right) \\
& \times\left(\log\left(2 M_{1} M_{2}\right)\right)^{2} .
\end{aligned}
$$

\(\blacksquare\)

\section*{4. The Spacing Problem}

In this section we investigate the distribution of real numbers of type
\[
  t(m,q)=(m+q)^\alpha-(m-q)^\alpha
\]
with $\alpha\ne0,1$, where $m,q$ range over integers with
$M\le m<2M$, $Q\le q<2Q$, and $3Q<M$. Note that
\[
  |t(m,q)|\asymp M^{\alpha-1}Q=T,
\]
say. Let $\mathscr B(M,Q,\Delta)$ be the number of quadruples
$(m,\widetilde m,q,\widetilde q)$ such that
\begin{equation*}
  |t(m,q)-t(\widetilde m,\widetilde q)|<\Delta T.
  \tag{4.1}
\end{equation*}
Our aim is to prove the following.

\medskip
\noindent\textbf{Proposition 2.}
If $Q\le M^{2/3}$, then
\begin{equation*}
  \mathscr B(M,Q,\Delta)
  \ll\left(MQ+\Delta M^2Q^2+M^{-2}Q^6\right)(\log(2M))^4,
  \tag{4.2}
\end{equation*}
where the constant implied by $\ll$ depends on $\alpha$ only.

The proof depends on three lemmas.

\medskip
\noindent\textbf{Lemma 3.}
Let $\mathscr C(A,B,M,\Delta)$ be the number of integers $m$ with
$M\le m<2M$ such that
\[
  \|Am-Bm^{-1}\|<\Delta.
\]
If $0<B<\Delta M^2$, then
\[
  \begin{aligned}
  \mathscr C(A,B,M,\Delta)\ll{}&
  \Delta M\sum_{0\le s<\Delta^{-1}}(1+\|sA\|M)^{-1}\\
  &+\Delta B^{-1/2}M^{3/2}
  \sum_{\substack{0<s<\Delta^{-1}\\
                   \|sA\|<2sBM^{-2}}}s^{-1/2},
  \end{aligned}
\]
where the constant implied by $\ll$ is absolute.

\medskip
\noindent\textit{Proof.}
We may assume that $\Delta\le\frac14$, since otherwise the assertion is
trivial. For $S>0$ we have the identity
\[
  \sum_{|s|<S}\left(1-\frac{|s|}{S}\right)e(sx)
  =\frac{1-\{S\}}{S}\left(\frac{\sin(\pi x[S])}{\sin(\pi x)}\right)^2
  +\frac{\{S\}}{S}\left(\frac{\sin(\pi x[S+1])}{\sin(\pi x)}\right)^2.
\]
Hence the sum is positive for all real $x$, and it is at least $S/4$ if
$\|x\|<(4S)^{-1}$. It follows that
\[
  \mathscr C(A,B,M,\Delta)
  \ll S^{-1}\sum_{0\le s<S}
  \left|\sum_{M\le m<2M}e(Asm-Bsm^{-1})\right|,
\]
with $S=(4\Delta)^{-1}$. The constant term $s=0$ contributes
$O(\Delta M)$. For $1\le s\le S$ the innermost sum is equal to
\[
  \int_M^{2M}e(\pm\|As\|\xi-Bs\xi^{-1})\,d\xi+O(1)
\]
(see Lemma 4.8 in~[6]). By Lemma 4.4 of~[6],
\[
  \left|\int_M^{2M}e(\pm\|As\|\xi-Bs\xi^{-1})\,d\xi\right|
  \ll(sB)^{-1/2}M^{3/2},
\]
and by Lemma 4.2 of~[6],
\[
  \left|\int_M^{2M}e(\pm\|As\|\xi-Bs\xi^{-1})\,d\xi\right|
  \ll\|sA\|^{-1},
\]
unless $\|sA\|<2sBM^{-2}$. Finally, we have the trivial bound
\[
  \left|\int_M^{2M}e(\pm\|As\|\xi-Bs\xi^{-1})\,d\xi\right|\le M.
\]
Gathering these estimates completes the proof. \(\blacksquare\)

\medskip
\noindent\textbf{Lemma 4.}
Let
\[
  A(q/\widetilde q)=(q/\widetilde q)^\gamma,\qquad \gamma\ne0,
\]
and let $B(q,\widetilde q)\asymp|q-\widetilde q|Q$ for
$Q\le q,\widetilde q<2Q$. Let $\mathscr D(M,Q,\Delta)$ be the number of
triples $(m,q,\widetilde q)$ of integers in the ranges
$M\le m<2M$, $Q\le q,\widetilde q<2Q$ such that
\[
  \|A(q/\widetilde q)m-B(q,\widetilde q)m^{-1}\|<\Delta.
\]
If $Q<M^{2/3}$, then
\[
  \mathscr D(M,Q,\Delta)
  \ll\left(MQ+\Delta MQ^2+Q^{8/3}\right)(\log(2M))^4.
\]

\medskip
\noindent\textit{Proof.}
Without loss of generality we may assume that $Q^{-1}<\Delta<1$. First we
count the triples $(m,q,\widetilde q)$ with
$|B(q,\widetilde q)|<\Delta M$ by a crude argument. We have
\[
  \|A(q/\widetilde q)m\|<2\Delta,
\]
which implies that there is an integer $\widetilde m$ such that
\[
  \left|\left(\frac q{\widetilde q}\right)^\gamma
        -\frac{\widetilde m}{m}\right|\ll\Delta M^{-1}.
\]
Therefore, by Lemma~1, the number of such triples is
\[
  O\left(MQ\log(2M)+\Delta MQ^2\right).
\]
For the remaining triples, those with
$|B(q,\widetilde q)|\ge\Delta M$, Lemma~3 gives
\[
  O\left(\Delta MQ^2+\mathscr D_1(M,Q,\Delta)
                  +\mathscr D_2(M,Q,\Delta)\right),
\]
where
\[
  \mathscr D_1(M,Q,\Delta)
  =\Delta M\sum_{0<s<\Delta^{-1}}
    \sum_{Q\le q,\widetilde q<2Q}
    \bigl(1+\|sA(q/\widetilde q)\|M\bigr)^{-1},
\]
and
\[
  \mathscr D_2(M,Q,\Delta)
  =\Delta M^{3/2}Q^{-1/2}
  \sum_{0<s<\Delta^{-1}}^{*}
  \sum_{Q\le q<2Q}^{*}\sum_{Q\le\widetilde q<2Q}^{*}
  (s|q-\widetilde q|)^{-1/2}.
\]
Here the stars mean that the summation is restricted by
\[
  \Delta M\ll |q-\widetilde q|Q
  \qquad\text{and}\qquad
  \|sA(q/\widetilde q)\|\ll s|q-\widetilde q|QM^{-2}.
\]
First we estimate $\mathscr D_1$. Split the range of summation into
subsets defined by
\[
  S\le s<2S,
  \qquad
  \|sA(q/\widetilde q)\|<\delta,
\]
where $1\le S<\Delta^{-1}$ and $M^{-1}<\delta<1$. The complete splitting
can be arranged with at most $O((\log(2M))^2)$ subsets, so
\[
  \mathscr D_1(M,Q,\Delta)
  \ll \Delta\delta^{-1}\mathscr A(Q,S,\delta S^{-1})
       (\log(2M))^2
\]
for some relevant $S$ and $\delta$. Since $Q\le M$, Lemma~1 yields
\[
  \mathscr D_1(M,Q,\Delta)\ll MQ(\log(2M))^3.
\]
Now estimate $\mathscr D_2$ similarly, splitting the range into subsets
defined by
\[
  S\le s<2S,
  \qquad
  R\le|q-\widetilde q|<2R,
\]
where $1\le S<\Delta^{-1}$ and
$\Delta MQ^{-1}\ll R\le Q$. We obtain
\[
  \mathscr D_2(M,Q,\Delta)
  \ll \Delta M^{3/2}(QRS)^{-1/2}
       \mathscr A(Q,S,M^{-2}QR)(\log(2M))^2
\]
for some relevant $S$ and $R$. By Lemma~1,
\[
  \mathscr D_2(M,Q,\Delta)
  \ll\left(MQ+(\Delta M)^{-1/2}Q^3\right)(\log(2M))^4.
\]
Hence
\[
  \mathscr D(M,Q,\Delta)
  \ll \Delta MQ^2+
  \left(MQ+(\Delta M)^{-1/2}Q^3\right)(\log(2M))^4.
\]
Finally, since $\mathscr D(M,Q,\Delta)$ is nondecreasing in $\Delta$, we
may replace $\Delta$ on the right-hand side by
$\Delta+M^{-1}Q^{2/3}$, completing the proof. \(\blacksquare\)

We are now ready to prove Proposition~2. Clearly (4.1) implies
\[
  |s(m,q)-s(\widetilde m,\widetilde q)|
  \ll \Delta T^{1/(\alpha-1)}
  =\Delta M Q^{1/(\alpha-1)},
\]
where
\[
  s(m,q)=\left(\frac{t(m,q)}{2\alpha}\right)^{1/(\alpha-1)}
  =q^{1/(\alpha-1)}m f(q/m),
\]
with
\[
  f(u)=\left(\frac{(1+u)^\alpha-(1-u)^\alpha}{2\alpha u}\right)^{1/(\alpha-1)}
  =1+\frac{\alpha-2}{6}u^2+O(u^4).
\]
Hence
\[
  \begin{aligned}
  \Bigg|{}&q^{1/(\alpha-1)}m
  -\widetilde q^{1/(\alpha-1)}\widetilde m\\
  &+\frac{\alpha-2}{6}
    \left(q^{(2\alpha-1)/(\alpha-1)}m^{-1}
    -\widetilde q^{(2\alpha-1)/(\alpha-1)}\widetilde m^{-1}\right)
  \Bigg|\\
  &\ll \Delta M Q^{1/(\alpha-1)}
       +M^{-3}Q^{4+1/(\alpha-1)}.
  \end{aligned}
\]
This gives the first approximation to $\widetilde m$ in terms of
$m,q,\widetilde q$, namely
\[
  \widetilde m
  =\left(\frac q{\widetilde q}\right)^{1/(\alpha-1)}m
   +O(\Delta M+M^{-1}Q^2),
\]
which, inserted into the last term, yields the second, stronger
approximation
\[
  |A(q/\widetilde q)m-\widetilde m+B(q,\widetilde q)m^{-1}|
  \ll\Delta M+M^{-3}Q^4,
\]
where
\[
  A(q/\widetilde q)=(q/\widetilde q)^{1/(\alpha-1)}
\]
and
\[
  B(q,\widetilde q)=\frac{\alpha-2}{6}
  \left(q^2A(q/\widetilde q)
       -\widetilde q^2A(\widetilde q/q)\right).
\]
The second approximation implies
\begin{equation*}
  \|A(q/\widetilde q)m+B(q,\widetilde q)m^{-1}\|
  \ll\Delta M+M^{-3}Q^4,
  \tag{4.3}
\end{equation*}
and, for given $m,q,\widetilde q$, the number of possible
$\widetilde m$ is $O(1+\Delta M)$. By Lemma~4,
\[
  \mathscr B(M,Q,\Delta)
  \ll(1+\Delta M)
  \left(MQ+\Delta M^2Q^2+M^{-2}Q^6+Q^{8/3}\right)(\log(2M))^4.
\]
Since
\[
  Q^{8/3}=(MQ)^{2/3}(M^{-2}Q^6)^{1/3}
  \le MQ+M^{-2}Q^6,
\]
the last term may be omitted. If $\Delta M<1$, we obtain (4.2);
otherwise the trivial bound gives
\[
  \mathscr B(M,Q,\Delta)
  \ll(1+\Delta M)MQ^2\ll\Delta M^2Q^2.
\]
This completes the proof of Proposition~2. \(\blacksquare\)

\section*{5. A Combination of Weyl's Shift with Poisson's Summation}
In this section we anticipate the use of Lemma~2 by an application of
Poisson's summation and the stationary phase method to the sum
\[
S_{\varphi\psi}(M_1,M_2,M_3,M_4)
 =\sum_{\substack{m_j\sim M_j\\1\le j\le4}}
   \varphi_{m_1m_2}\psi_{m_3}
   e\left(x\frac{m_1^{\alpha_1}m_2^{\alpha_2}m_3^\alpha m_4^{-\alpha}}
                  {M_1^{\alpha_1}M_2^{\alpha_2}M_3^\alpha M_4^{-\alpha}}\right).
\]
We appeal to a special case of the van der Corput lemma.

\medskip
\noindent\textbf{Lemma 5.}
Let $X>0$, $M>0$, $\mu>1$, and $\alpha\ne0,1$. We then have
\[
\begin{aligned}
\sum_{M<m<\mu M}m^{-1/2}
 e\left(\alpha^{-1}m^\alpha M^{-\alpha}X\right)
={}&\gamma\sum_{N<n<\nu N}n^{-1/2}
 e\left(-\beta^{-1}n^\beta N^{-\beta}X\right)\\
&+O\left(M^{-1/2}\log(2+M)+N^{-1/2}\log(2+N)\right),
\end{aligned}
\]
with $\beta=\alpha/(\alpha-1)$, $\nu=\mu^{\alpha-1}$, $MN=X$,
and some $\gamma$ depending on $\alpha$ alone. The constant implied in the
symbol $O$ depends at most on $\alpha$ and $\mu$. If $\nu<1$
($\alpha<1$), the range of summation in $n$ is understood to be
$\nu N<n<N$.

\medskip
\noindent\textit{Proof.}
This result is a special case of Theorem~4.9 of [6], except that we claim a
better error term, for which see [7]. \(\blacksquare\)

In our applications of Lemma~5 the parameter $X$ will depend
multiplicatively on several variables; so will $N$. In order to separate the
dependence from the range of summation we appeal to the following formula.

\medskip
\noindent\textbf{Lemma 6.}
Let $0<L\le N<\nu N<\lambda L$ and let $a_l$ be complex numbers with
$|a_l|\le1$. We then have
\[
\sum_{N<n<\nu N}a_n
 =\frac1{2\pi}\int_{-L}^{L}
   \left(\sum_{L<l<\lambda L}a_l l^{-it}\right)
   N^{it}(\nu^{it}-1)t^{-1}\,dt+O(\log(2+L)),
\]
where the constant implied in $O$ depends on $\lambda$ only.
\(\blacksquare\)

Combining Lemmas~5 and~6 we obtain the following.

\medskip
\noindent\textbf{Lemma 7.}
Let $X>0$, $M>0$, $\mu>1$, and $\alpha\ne0,1$. We then have
\[
\begin{aligned}
\sum_{M<m<\mu M}m^{-1/2}
 e\left(\alpha^{-1}m^\alpha M^{-\alpha}X\right)
={}&\frac{\gamma}{2\pi}\int_{-L}^{L}
 \left(\sum_{\lambda^{-1}L<l<\lambda L}
 l^{-1/2}\left(\frac{X}{lM}\right)^{it}
 e\left(-\beta^{-1}\left(\frac{lM}{X}\right)^\beta X\right)\right)\\
&\hspace{4.5em}\times
 \frac{\mu^{i(\alpha-1)t}-1}{t}\,dt\\
&+O\left(M^{-1/2}\log(2+M)+L^{-1/2}\log(2+L)\right),
\end{aligned}
\]
where $\beta=\alpha/(\alpha-1)$,
$\lambda=2(\mu^{\alpha-1}+\mu^{1-\alpha})$, and $L$ is any number with
$\frac12<LMX^{-1}<2$. The constant implied in the symbol $O$ depends at
most on $\alpha$ and $\mu$.

Now, by Lemma~7 with $\alpha$ replaced by $-\alpha$, we obtain
\[
\begin{aligned}
S_{\varphi\psi}\ll{}&x^{-1/2}M_4
 \sum_{m_1\sim M_1}\sum_{m_2\sim M_2}\sum_{m_3\sim M_3}\sum_{l\sim L}
 \widetilde\varphi_{m_1m_2}\widetilde\psi_{m_3}\zeta_l\\
&\quad\times e\left(\widetilde\alpha^{-1}x
 \frac{m_1^{\beta_1}m_2^{\beta_2}m_3^\beta l^\beta}
      {M_1^{\beta_1}M_2^{\beta_2}M_3^\beta L^\beta}\right)\\
&+O\left(M_1M_2M_3(1+x^{-1/2}M_4)\log(xM_4)\right),
\end{aligned}
\]
where $\widetilde\varphi_{m_1m_2}$ and $\widetilde\psi_{m_3}$ are some
complex numbers with $|\widetilde\varphi_{m_1m_2}|\le1$ and
$|\widetilde\psi_{m_3}|\ll\log(2x)$, $\zeta_l$ is a complex number with
$|\zeta_l|\le1$, and
\[
 \beta_1=\frac{\alpha_1}{\alpha+1},\qquad
 \beta_2=\frac{\alpha_2}{\alpha+1},\qquad
 \beta=\frac{\alpha}{\alpha+1},\qquad
 \widetilde\alpha=(\alpha+1)\alpha^\beta,\qquad
 L=M_4^{-1}x.
\]
Since the exponents in $m_3$ and $l$ are equal, this permits us to treat
$l m_3$ as one variable, $m=l m_3$, say, with the multiplicity bounded by
the divisor function $\tau(m)\ll m^\varepsilon$. Theorem~3 is applicable
with $M=LM_3=xM_3M_4^{-1}$, giving the following.

\medskip
\noindent\textbf{Theorem 5.}
Let $\alpha,\alpha_1,\alpha_2$ be real constants such that $\alpha\ne-1$
and $\alpha\alpha_1\alpha_2\ne0$. Let $M_1,M_2,M_3,M_4,x\ge1$, and let
$\varphi_{m_1m_2},\psi_{m_3}$ be complex numbers with
$|\varphi_{m_1m_2}|\le1$ and $|\psi_{m_3}|\le1$. We then have
\[
\begin{aligned}
S_{\varphi\psi}(M_1,M_2,M_3,M_4)\ll{}&
 \left\{x^{1/2}(M_1M_2)^{3/4}M_3
 +x^{7/20}M_1M_2M_3^{11/10}M_4^{-1/10}\right.\\
&\quad+x^{1/4}(M_1M_2)^{3/4}(M_3M_4)^{1/2}
 +x^{1/5}M_1M_2M_3^{7/10}M_4^{3/10}\\
&\quad\left.+x^{-1/2}M_1M_2M_3M_4\right\}
 (xM_1M_2M_3M_4)^\varepsilon
\end{aligned}
\]
for any $\varepsilon>0$; the constant implied in $\ll$ depends on
$\alpha,\alpha_1,\alpha_2$, and $\varepsilon$ only.

\section*{6. A Special Sum}
In this section we estimate sums of type
\[
S_{\chi\varphi\psi}(H,M,N)
 =\sum_{h\sim H}\sum_{m\sim M}\sum_{n\sim N}
  \chi(h)\varphi_m\psi_n
  e\left(x\frac{h n^{-1}m^\alpha}{H N^{-1}M^\alpha}\right),
\]
where $\chi(h)$ is an additive character, i.e., $\chi(h)=e(\xi h)$, say.
Such sums occur in applications of modern sieve methods. The nature of the
monomial $h n^{-1}$ will be exploited effectively by a refined version of
the argument established in [5]. We prove the following.

\medskip
\noindent\textbf{Theorem 6.}
Let $\alpha\ne0,1$ and $H,M,N,x\ge1$. Let $\chi(h)$ be an additive
character and let $\varphi_m,\psi_n$ be complex numbers with
$|\varphi_m|\le1$ and $|\psi_n|\le1$. We then have
\[
\begin{aligned}
S_{\chi\varphi\psi}(H,M,N)\ll{}&(HMN)^{1/2}
 \Bigl[(H+N)^{1/2}
 \bigl(x^{1/8}H^{-1/6}M^{1/12}N^{1/6}
      +x^{1/8}H^{-1/8}N^{3/8}\\
&\hspace{10em}{}+N^{1/2}+N^{1/4}M^{1/8}\bigr)x^{1/8}
 +M^{1/2}+x^{-1/4}M^{1/2}N\Bigr]\mathscr L^4,
\end{aligned}
\]
where the constant implied in $\ll$ depends on $\alpha$ only, and we set
$\mathscr L=\log(2HMNx)$.

The proof is rather long. We begin by applying Cauchy's inequality:
\[
\begin{aligned}
|S_{\chi\varphi\psi}|^2\ll{}&M^{3/2}\sum_{m\sim M}m^{-1/2}
 \left|\sum_{h\sim H}\sum_{n\sim N}\chi(h)\psi_n
 e\left(x\frac{h n^{-1}m^\alpha}{H N^{-1}M^\alpha}\right)\right|^2\\
={}&M^{3/2}\sum_{n_1\sim N}\sum_{n_2\sim N}
 \psi_{n_1}\overline{\psi_{n_2}}\sum_k\omega_{n_1n_2}(k)
 \sum_{m\sim M}m^{-1/2}e\left(y\frac{k m^\alpha}{n_1n_2}\right),
\end{aligned}
\]
say, where $y=xH^{-1}NM^{-\alpha}$ and $\omega_{n_1n_2}(k)$ is defined by
\[
\omega_{n_1n_2}(k)
 =\sum_{\substack{h_1\sim H,\ h_2\sim H\\h_1n_2-h_2n_1=k}}
  \chi(h_1)\overline{\chi(h_2)}.
\]
The terms with $k=0$ contribute trivially $O(HM^2N\mathscr L)$. The
remaining $k$ fall into dyadic intervals $|k|\sim K$ with
$1\ll K\ll HN$. For technical convenience we wish to work with $n_1,n_2$
coprime. To make that restriction, observe that
$\omega_{n_1n_2}(k)=\omega_{n_1^*n_2^*}(k^*)$, where
$k^*=k/d$, $n_1^*=n_1/d$, and $n_2^*=n_2/d$, with $d=(n_1,n_2)$.
We get
\[
\begin{aligned}
|S_{\chi\varphi\psi}|^2\ll{}&\mathscr L^3M^{3/2}
 \sum_{d\sim D}\sum_{n_1\sim N_1}
 \sum_{\substack{n_2\sim N_2\\(n_1,n_2)=1}}
 \left|\sum_{r\sim R}\omega_{n_1n_2}(r)
 \sum_{m\sim M}m^{-1/2}
 e\left(y\frac{r m^\alpha}{d n_1n_2}\right)\right|\\
&+HM^2N\mathscr L^3
\end{aligned}
\]
for some $K$ with $1\ll K\ll HN$ and for some $D$ with
$1\le D\le\min\{K,N\}$, where $R=D^{-1}K$ and
$N_1=N_2=D^{-1}N$. Hence, by Lemma~7, we obtain
\begin{equation*}
\begin{aligned}
|S_{\chi\varphi\psi}|^2\ll{}&\mathscr L^3M^2
 \left(\frac{HN}{xK}\right)^{1/2}
 \sum_{d\sim D}\sum_{\substack{n_1\sim N_1,\ n_2\sim N_2\\(n_1,n_2)=1}}
 \sum_{l\sim L}
 \left|\sum_{r\sim R}r^{it}\omega_{n_1n_2}(r)
 e\left(wr^{-\beta}R^\beta\right)\right|\\
&+\mathscr L^3HMN\left(M+N^2+x^{-1/2}MN^2\right),
\end{aligned}
\tag{6.1}
\end{equation*}
where $L=|\alpha|xK(HMN)^{-1}$, $t$ is a real number,
$\beta=(\alpha-1)^{-1}$, and
\[
w=(1-\alpha)\frac{xK}{HN}
 \left(\frac lL\right)^{\alpha\beta}
 \left(\frac{dn_1n_2}{DN_1N_2}\right)^\beta.
\]

Our nearest aim is now to build a single variable $c=dn_1n_2$, say, out
of the three variables $d,n_1,n_2$. We begin by Fourier analysis of the
arithmetic function $\omega_{n_1n_2}(r)$, separating the variable $r$ from
the parameters $n_1,n_2$.

\medskip
\noindent\textbf{Lemma 8.}
Let $H_1\ge H_1'\ge1$, $H_2\ge H_2'\ge1$, let $\xi_1,\xi_2$ be real
numbers, and let $n_1,n_2$ be positive integers with $(n_1,n_2)=1$.
We have
\begin{equation*}
\begin{aligned}
\omega(r)
 &=\sum_{\substack{H_1'\le h_1<H_1,\ H_2'\le h_2<H_2\\
                    h_1n_1-h_2n_2=r}}
   e(\xi_1h_1+\xi_2h_2)\\
 &=\int_0^1\widehat\omega(\theta)e(\theta r)\,d\theta,
\end{aligned}
\tag{6.2}
\end{equation*}
with
\begin{equation*}
 \int_0^1|\widehat\omega(\theta)|\,d\theta
 \ll\left(1+\frac{H_1H_2}{n_1n_2}\right)^{1/2}
       \bigl(\log(2H_1H_2)\bigr)^2,
\tag{6.3}
\end{equation*}
and the constant implied in $\ll$ is absolute.

\medskip
\noindent\textit{Proof.}
We have (6.2) with
\[
\widehat\omega(\theta)
 =\left(\sum_{H_1'\le h_1<H_1}e((\xi_1-\theta n_1)h_1)\right)
  \left(\sum_{H_2'\le h_2<H_2}e((\xi_2+\theta n_2)h_2)\right).
\]
Since
\[
\left|\sum_{M_1\le m<M}e(\xi m)\right|
 =\sum_{m=-\infty}^{\infty}\alpha(m)e(\xi m),
\]
with
\[
 \alpha(m)=\alpha(M_1,M;m)
 \ll\left(1+\frac{m^2}{M^2}\right)^{-1}\log(2M),
\]
it gives
\[
|\widehat\omega(\theta)|
 =\left(\sum_{h_1=-\infty}^{\infty}\alpha_1(h_1)
        e((\xi_1-\theta n_1)h_1)\right)
  \left(\sum_{h_2=-\infty}^{\infty}\alpha_2(h_2)
        e((\xi_2+\theta n_2)h_2)\right),
\]
with
\[
 \alpha_j(h_j)\ll(1+h_j^2H_j^{-2})^{-1}\log(2H_j).
\]
Hence
\[
\begin{aligned}
\int_0^1|\widehat\omega(\theta)|\,d\theta
={}&\sum_{\substack{h_1,h_2\in\mathbb Z\\h_1n_1=h_2n_2}}
 \alpha_1(h_1)\alpha_2(h_2)e(\xi_1h_1+\xi_2h_2)\\
\ll{}&\sum_{h=-\infty}^{\infty}
 (1+h^2n_2^2H_1^{-2})^{-1}
 (1+h^2n_1^2H_2^{-2})^{-1}
 \bigl(\log(2H_1H_2)\bigr)^2\\
\ll{}&\left(1+\min\left\{\frac{H_1}{n_2},\frac{H_2}{n_1}\right\}\right)
 \bigl(\log(2H_1H_2)\bigr)^2,
\end{aligned}
\]
and this yields (6.3). \(\blacksquare\)

By Lemma~8 and the Cauchy--Schwarz inequality we deduce that
\[
\begin{aligned}
\left|\sum_{r\sim R}r^{it}\omega_{n_1n_2}(r)
 e(wr^{-\beta}R^\beta)\right|^2
\le{}&\left(\int_0^1|\widehat\omega(\theta)|\,d\theta\right)\\
&\times\int_0^1|\widehat\omega(\theta)|
 \left|\sum_{r\sim R}r^{it}
 e(\theta r+wr^{-\beta}R^\beta)\right|^2d\theta.
\end{aligned}
\]
Next, by Lemma~2, we obtain
\[
\begin{aligned}
\left|\sum_{r\sim R}r^{it}e(\theta r+wr^{-\beta}R^\beta)\right|^2
\ll{}&\frac{R^2}{Q}+\frac{R}{Q}\sum_{1\le q<Q}
 \left|\sum_{\substack{r+q\sim R\\r-q\sim R}}
 \left(\frac{r+q}{r-q}\right)^{it}
 e\left(wR^\beta\bigl[(r+q)^{-\beta}-(r-q)^{-\beta}\bigr]\right)\right|
\end{aligned}
\]
with any $Q<\frac13R$. Note that the bound does not depend on $\theta$ and
that
\[
 \int_0^1|\widehat\omega(\theta)|\,d\theta
 \ll(1+DHN^{-1})\mathscr L^2.
\]
Combining the above results with (6.1), we conclude that
\[
\begin{aligned}
|S_{\chi\varphi\psi}|^2\ll{}&\mathscr L^5M^2
 \left(\frac{HN}{xK}\right)^{1/2}\left(1+\frac{DH}{N}\right)
 \sum_{d\sim D}\sum_{n_1\sim N_1}\sum_{n_2\sim N_2}\sum_{l\sim L}\\
&\times\left(\frac{R^2}{Q}+\frac{R}{Q}\sum_{1\le q<Q}
 \left|\sum_{\substack{r+q\sim R\\r-q\sim R}}
 \left(\frac{r+q}{r-q}\right)^{it}
 e\left(wR^\beta\bigl[(r+q)^{-\beta}-(r-q)^{-\beta}\bigr]\right)
 \right|\right)^{1/2}\\
&+\mathscr L^3HMN\left(M+N^2+x^{-1/2}MN^2\right).
\end{aligned}
\]
Note that the variables $d,n_1,n_2$ occur only in $w$ as a product
$c=dn_1n_2$. Given $c\sim DN_1N_2=D^{-1}N^2=C$, say, the number of
representations of $c$ in the form $c=dn_1n_2$ is bounded by the divisor
function $\tau_3(c)$. By Cauchy's inequality we obtain
\[
\begin{aligned}
|S_{\chi\varphi\psi}|^4\ll{}&\mathscr L^{12}D^{-1}M^3(DH+N)^2
 \left(\frac{xNK^3}{HMQD^3}+\frac{K}{QD}\sum_{1\le q<Q}S_q\right)\\
&+\mathscr L^6(HMN)^2\left(M+N^2+x^{-1/2}MN^2\right)^2,
\end{aligned}
\]
where
\[
S_q=\sum_{c\sim C}\sum_{l\sim L}
 \left|\sum_{\substack{r+q\sim R\\r-q\sim R}}
 \left(\frac{r+q}{r-q}\right)^{it}
 e\left(wR^\beta\bigl[(r+q)^{-\beta}-(r-q)^{-\beta}\bigr]\right)\right|,
\]
with
\[
 w=(1-\alpha)\frac{xK}{HN}
   \left(\frac lL\right)^{\alpha\beta}
   \left(\frac cC\right)^\beta,
 \qquad \beta=(\alpha-1)^{-1}.
\]
It remains to estimate $S_q$ for each $q$ individually. We appeal to
Proposition~1. The arguments are similar to those used in the proof of
Theorem~2 if we treat
\[
 \delta(r)=(r+q)^{-\beta}-(r-q)^{-\beta}\asymp q r^{-\alpha\beta}
\]
as a monomial in $r$ of degree $-\alpha\beta$. More to the point, we have
\[
 |\delta(r)-\delta(\widetilde r)|
 \asymp q\left|r^{-\alpha\beta}-\widetilde r^{-\alpha\beta}\right|,
\]
so the spacing problems for the points $\delta(r)$ and for the points
$r^{\alpha\beta}$ are equivalent. From these remarks it is clear that
Proposition~1 produces a bound for $S_q$ of the same form as that in
Theorem~2, with the parameters $M_1=L$, $M_2=C$, $M_3=R$, $M_4=1$, and
with $x$ replaced by
\[
 \widetilde x=qxDH^{-1}N^{-1}.
\]
We obtain
\[
\begin{aligned}
\mathscr L^{-1}S_q\ll{}&
 (\widetilde xM_1M_2M_3M_4)^{1/2}
 +M_1M_2(M_3M_4)^{1/2}
 +(M_1M_2)^{1/2}M_3M_4\\
&+\widetilde x^{-1/2}M_1M_2M_3M_4\\
\ll{}&\frac{xK}{H}\left(\frac{q}{DM}\right)^{1/2}
 +\frac{xN}{HM}\left(\frac KD\right)^{3/2}
 +\left(\frac{xN}{HM}\right)^{1/2}\left(\frac KD\right)^{3/2}\\
&+\left(\frac{x}{qH}\right)^{1/2}
  \frac{K^2N^{3/2}}{MD^{5/2}}.
\end{aligned}
\]
Hence
\[
\begin{aligned}
|S_{\chi\varphi\psi}|^4\ll{}&\mathscr L^{13}D^{-1}M^3(DH+N)^2
 \left(W_1Q^{-1}+W_2Q^{-1/2}+W_3+W_4Q^{1/2}\right)\\
&+\mathscr L^6(HMN)^2\left(M+N^2+x^{-1/2}MN^2\right)^2,
\end{aligned}
\]
where
\[
\begin{aligned}
W_1&=\frac{xN}{HM}\left(\frac KD\right)^3,\\
W_2&=\left(\frac{xN^3K^6}{D^7HM^2}\right)^{1/2},\\
W_3&=\frac{xN}{HM}\left(\frac KD\right)^{5/2}
 +\left(\frac{xN}{HM}\right)^{1/2}\left(\frac KD\right)^{5/2},\\
W_4&=\frac{xK^2}{HM^{1/2}D^{3/2}}.
\end{aligned}
\]
We choose
\[
 Q=\min\left\{\frac{K}{3D},\frac{W_2}{W_4}
       +\left(\frac{W_1}{W_4}\right)^{2/3}\right\},
\]
giving
\[
\begin{aligned}
W_1Q^{-1}+W_2Q^{-1/2}+W_3+W_4Q^{1/2}
\ll{}&W_1\frac DK+W_2\left(\frac DK\right)^{1/2}+W_3
 +W_1^{1/3}W_4^{2/3}+W_2^{1/2}W_4^{1/2}\\
\ll{}&\frac{xNK^2}{HMD^2}
 +\left(\frac{xN^3K^5}{HM^2D^6}\right)^{1/2}
 +\frac{xN}{HM}\left(\frac KD\right)^{5/2}\\
&+\left(\frac{xN}{HM}\right)^{1/2}\left(\frac KD\right)^{5/2}
 +\frac{x}{H}\left(\frac{NK^7}{M^2D^6}\right)^{1/3}\\
&+\left(\frac{xN}{HM}\right)^{3/4}\left(\frac KD\right)^{5/2}.
\end{aligned}
\]
Here the first term is majorized by the third term, and the sixth term is
equal to the geometric mean of the third and fourth terms. Thus we obtain
\[
\begin{aligned}
|S_{\chi\varphi\psi}|^4\ll{}&\mathscr L^{13}D^{-1}M^3(DH+N)^2
 \left\{\left(\frac{xN^3K^5}{HM^2D^6}\right)^{1/2}
 +\frac{xN}{HM}\left(\frac KD\right)^{5/2}\right.\\
&\left.\hspace{7em}
 +\left(\frac{xN}{HM}\right)^{1/2}\left(\frac KD\right)^{5/2}
 +\frac{x}{H}\left(\frac{NK^7}{M^2D^6}\right)^{1/3}\right\}\\
&+\mathscr L^6(HMN)^2\left(M+N^2+x^{-1/2}MN^2\right)^2.
\end{aligned}
\]
Now observe that the worst case is $D\asymp1$ and $K\asymp HN$, giving
\[
\begin{aligned}
|S_{\chi\varphi\psi}|^4\ll{}&\mathscr L^{13}M^2(H+N)^2
 \left\{x^{1/2}H^2N^4+xH^{3/2}N^{7/2}
 +x^{1/2}H^2M^{1/2}N^3\right.\\
&\left.\hspace{10em}{}+xH^{4/3}M^{1/3}N^{8/3}\right\}\\
&+\mathscr L^6(HMN)^2\left(M+N^2+x^{-1/2}MN^2\right)^2.
\end{aligned}
\]
Here the last but one term is majorized by the first term, so it can be
omitted. This completes the proof of Theorem~6. \(\blacksquare\)

\section*{7. Exponential Sums Related to Short Intervals}
Considerable attention in analytic number theory is given to the
distribution of special sequences in short intervals
\[
 \mathscr A=\{n:x-y<n\le x\},
\]
with $1\le y\le\frac12x$. A powerful approach to many such problems is
offered by the sieve method, which depends on estimates for the remainder
terms
\[
 r_d=\left[\frac xd\right]-\left[\frac{x-y}{d}\right]-\frac yd.
\]
The modern version of the linear sieve [4] requires estimates for the
bilinear form
\begin{equation*}
 R(M,N)=\sum_{1\le m<M}\sum_{1\le n<N}a_mb_nr_{mn}.
\tag{7.1}
\end{equation*}
Here the real coefficients $a_m,b_n$ are subject to $|a_m|\le1$ and
$|b_n|\le1$. One needs the upper bound
\begin{equation*}
 R(M,N)\ll yx^{-\varepsilon}.
\tag{7.2}
\end{equation*}
Clearly (7.2) holds true if $MN\le yx^{-\varepsilon}$ because $|r_d|\le1$.
However, one can do better by taking into account cancellation among the
terms in (7.1) due to the variation in sign of $r_d$. The following lemma
transforms the problem to estimating exponential sums of type
\[
 S_z(H,M,N)=\sum_{1\le h\le H}\sum_{1\le m<M}\sum_{1\le n<N}
 \frac{a_mb_n}{mn}e\left(\frac{hz}{mn}\right).
\]

\medskip
\noindent\textbf{Lemma 9.}
Let $M,N\ge1$ and $|a_m|\le1$, $|b_n|\le1$. We have
\[
 |R(M,N)|\le2\int_{x-2y}^{x+y}|S_z(H,M,N)|\,dz
 +O(yx^{-\varepsilon}),
\]
with $H=MNy^{-1}x^{3\varepsilon}$, the constant implied in $O$ depending
on $\varepsilon$ alone.

\medskip
\noindent\textit{Proof.}
Let $f(u)$ be a smooth function supported in the interval
\[
 [x-y-yx^{-2\varepsilon},\,x+yx^{-2\varepsilon}]
\]
such that $f(u)=1$ in $[x-y,x]$ and
\[
 f^{(j)}(u)\ll(yx^{-2\varepsilon})^{-j}
\]
for any $j\ge0$, the constant implied in $\ll$ depending on $j$ only.
We then have
\[
\begin{aligned}
\sum_{x-y<lmn\le x}a_mb_n
 &=\sum_l\sum_m\sum_n a_mb_nf(lmn)+O(yx^{-\varepsilon})\\
 &=\sum_m\sum_n\frac{a_mb_n}{mn}
   \sum_h\widehat f\left(\frac{h}{mn}\right)+O(yx^{-\varepsilon})
\end{aligned}
\]
by Poisson's summation, where $\widehat f(t)$ is the Fourier transform of
$f(u)$. The constant term ($h=0$) contributes
\[
 \widehat f(0)\sum_m\sum_n\frac{a_mb_n}{mn}
 =y\sum_m\sum_n\frac{a_mb_n}{mn}+O(yx^{-\varepsilon}).
\]
For $|h|>H=MNy^{-1}x^{3\varepsilon}$ we have
$\widehat f(h/(mn))\ll h^{-2}$, so the terms with $|h|>H$ contribute
$O(yx^{-\varepsilon})$. Combining these results completes the proof.
\(\blacksquare\)

By Theorem~6 and Lemma~9 we infer the following.

\medskip
\noindent\textbf{Theorem 7.}
Let $2\le y\le x^{1/2}$ and $|a_m|\le1$, $|b_n|\le1$. We then have
\begin{equation*}
 \sum_{M\le m<2M}\sum_{N\le n<2N}a_mb_nr_{mn}
 \ll yx^{-\varepsilon},
\tag{7.3}
\end{equation*}
provided
\begin{align*}
 M&<yx^{-\varepsilon'},                                      \tag{7.4}\\
 N^6&<My^7x^{-3-\varepsilon'},                               \tag{7.5}\\
 M^2N^4&<yx^{1-\varepsilon'},                                \tag{7.6}
\end{align*}
where $\varepsilon'=48\varepsilon$, the constant implied in $\ll$
depending on $\varepsilon$ alone.

\medskip
\noindent\textit{Proof.}
We have to show that
\[
 \sum_{h\sim H}\sum_{m\sim M}\sum_{n\sim N}
 a_mb_ne\left(\frac{hz}{mn}\right)\ll MNx^{-2\varepsilon}
\]
for any $H$ with $1\le H<MNy^{-1}x^{3\varepsilon}$ and
$x-2y<z<x+y$. Since $H<N$, we obtain by Theorem~6 that the sum is
\[
\begin{aligned}
\ll{}&(HMN)^{1/2}\Biggl[N^{1/2}
 \Biggl(\left(\frac{zH}{MN}\right)^{1/8}H^{-1/6}M^{1/12}N^{1/6}
 +\left(\frac{zH}{MN}\right)^{1/8}H^{-1/8}N^{3/8}\\
&\hspace{10em}{}+N^{1/2}+N^{1/4}M^{1/8}\Biggr)
 \left(\frac{zH}{MN}\right)^{1/8}
 +M^{1/2}+\left(\frac{MN}{zH}\right)^{1/4}M^{1/2}N\Biggr]
 (\log x)^4.
\end{aligned}
\]
Since $x-2y<z<x+y$ and $y\le x^{1/2}$, we have $z\asymp x$.
Now observe that the worst value of $H$ is
$H=MNy^{-1}x^{3\varepsilon}$, giving
\[
\begin{aligned}
\ll{}&y^{-1/2}MN\Bigl[x^{1/4}(yM)^{-1/12}N^{1/2}
 +x^{1/4}(yM)^{-1/8}N^{3/4}
 +x^{1/8}y^{-1/8}N\\
&\hspace{8em}{}+x^{1/8}y^{-1/8}N^{3/4}M^{1/8}
 +M^{1/2}+x^{-1/4}y^{1/4}M^{1/2}N\Bigr]x^{2\varepsilon}
 \ll MNx^{-2\varepsilon},
\end{aligned}
\]
provided (7.4), (7.5), and (7.6) hold, as well as
\begin{align*}
 N^6&<My^5x^{-2-\varepsilon'},                               \tag{7.7}\\
 N^8&<y^5x^{-1-\varepsilon'},                                \tag{7.8}\\
 MN^6&<y^5x^{-1-\varepsilon'}.                               \tag{7.9}
\end{align*}
But (7.7)--(7.9) follow from (7.4)--(7.6) under the hypothesis
$y\le x^{1/2}$. \(\blacksquare\)

The bilinear form (7.1) was investigated in various papers. It was proved
in [3,5] that (7.2) holds true subject to (7.4) and
\begin{equation*}
 M^2N^4<y^5x^{-1-\varepsilon'}.
\tag{7.10}
\end{equation*}
Combining (7.5) and (7.10), we conclude as follows.

\medskip
\noindent\textbf{Corollary.}
Let $x^{7/19}<y<x^{11/23}$ and $|a_m|\le1$, $|b_n|\le1$. We then have
(7.2), provided
\begin{equation*}
 M<yx^{-\varepsilon'}
\tag{7.4}
\end{equation*}
and
\begin{equation*}
 N<y^{19/16}x^{-7/16-\varepsilon'}.
\tag{7.11}
\end{equation*}

\section*{Acknowledgments}
The principal part of this work was accomplished in the summer of 1987,
when the first-named author visited Rutgers University. It is his pleasure
to express thanks to the Mathematics Department for the invitation and for
the financial support. The second author takes this opportunity to thank
Professor H. Halberstam and Professor H.-E. Richert for the stimulating
correspondence on related topics. Finally, the authors are grateful to the
referee for careful reading and for catching typographical errors in the
manuscript.

\section*{References}
\begin{itemize}
  \item[1.] E. Bombieri and H. Iwaniec, On the order of
  $\zeta\left(\frac12+it\right)$, \textit{Ann. Scuola Norm. Sup. Pisa}
  \textbf{13}, no.~3 (1986), 449--472.

  \item[2.] J. G. van der Corput, Verschärfung der Abschätzung beim
  Teilerproblem, \textit{Math. Ann.} \textbf{87} (1922), 39--65.

  \item[3.] H. Halberstam, D. R. Heath-Brown, and H.-E. Richert,
  Almost-primes in short intervals, in \textit{Recent Progress in Analytic
  Number Theory}, vol.~1, pp.~69--101, Academic Press, London, 1981.

  \item[4.] H. Iwaniec, A new form of the error term in the linear sieve,
  \textit{Acta Arith.} \textbf{37} (1980), 307--320.

  \item[5.] H. Iwaniec and M. Laborde, $P_2$ in short intervals,
  \textit{Ann. Inst. Fourier (Grenoble)} \textbf{31}, no.~4 (1981), 37--56.

  \item[6.] E. C. Titchmarsh, \textit{The Theory of the Riemann
  Zeta-Function}, Oxford, 1951.

  \item[7.] I. M. Vinogradov, \textit{Izbrannye Trudy}, Moscow, 1952;
  \textit{Selected Works}, Springer, Berlin, 1985.

  \item[8.] I. M. Vinogradov, A new method of estimation of trigonometrical
  sums, \textit{Mat. Sb.} \textbf{43} (1936), 175--188.

  \item[9.] H. Weyl, Über die Gleichverteilung von Zahlen mod.~Eins,
  \textit{Math. Ann.} \textbf{77} (1916), 313--352.
\end{itemize}

\end{document}
